\begin{solapartat}
\begin{minipage}[t]{0.7\linewidth}
\punts{0,5}

Les línies de camp magnètic són perpendiculars al pla de l'espira i sentit cap en fora del paper.

Cal justificar la direcció indicant que han de ser perpendiculars al corrent i al vector posició, i que el sentit es determina a partir de la regla de la mà dreta.
\end{minipage}\hfill
\begin{minipage}[t]{0.26\linewidth}
\vspace{0pt}
\figura[0.85]{sol_fig1.png}
\end{minipage}

\smallskip
\begin{minipage}[t]{0.7\linewidth}
També s'accepta un esquema del camp magnètic com el que apareix en el dibuix, on es representa el camp magnètic com a tangent a una circumferència centrada al fil i continguda al pla perpendicular a la direcció del fil. En qualsevol cas, s'ha d'indicar com s'obté el sentit del camp magnètic.
\end{minipage}\hfill
\begin{minipage}[t]{0.26\linewidth}
\vspace{0pt}
\figura[0.8]{sol_fig2.png}
\end{minipage}

\smallskip
Finalment, també es pot deduir a partir de: $d\vec{B} = \dfrac{\mu_0}{4\pi}\left[\dfrac{I d\vec{l}\times\vec{u}_r}{r^2}\right]$ indicant que tots els elements generen un camp en el mateix sentit i direcció.

\medskip
\punts{0,5} No, atès que el camp magnètic s'afebleix amb la distància. Com a justificació, n'hi ha prou a dir que la intensitat de camp magnètic disminueix quan ens allunyem de les fonts de camp.

Alternativament, també es pot donar la relació del camp creat per un fil infinit i indicar que és inversament proporcional a la distància:
\[ \vec{B} = \frac{\mu_0 I}{2\pi r} \]
\end{solapartat}

\begin{solapartat}
\punts{0,5} El camp magnètic és proporcional a la intensitat que circula pel fil, per tant, el flux de camp magnètic també és proporcional a la intensitat.

Perquè hi hagi un corrent induït cal que el flux de camp magnètic variï en el temps, per tant, quan el corrent és constant, també ho serà el flux i el corrent induït serà nul:
\begin{itemize}
  \item de 0 a 20s el corrent augmenta, hi haurà un corrent induït
  \item de 20 a 80~s el corrent roman constant, no hi haurà corrent induït
  \item de 80 a 120~s el corrent disminueix, hi haurà corrent induït.
\end{itemize}

\punts{0,5} La força electromotriu és la derivada temporal del flux, per tant, com més ràpid variï el corrent, més gran serà la derivada temporal i més intens el corrent induït. En el tram de 0 a 20s el corrent augmenta a un ritme constant de 0,05~A/s, mentre que en el darrer tram de 80 a 120 disminueix a un ritme de 0,025~A/s. Com que el ritme és més gran en el primer tram, el corrent induït també serà més intens en el tram de 0 a 20~s.

També es considerarà correcte si es determina correctament (i es justifica) el sentit del corrent induït en tots dos casos, i es dona com a correcte un dels resultats basant-se en el signe del corrent, és a dir, es dona com a valor més gran el positiu respecte el negatiu enlloc de basar-se en la magnitud del corrent.
\end{solapartat}
